% GENERATED FILE. Edit canonical lessons, then run npm run generate:books.
% canonical-source-hash: fnv1a64:8625b12f71c3f757
% canonical-lessons: ZH-C204-wangzhan, ZH-C204-denglu, ZH-C204-dianji, ZH-C204-shipin, ZH-C204-dazi

\chapter{Online}
\label{ch:zh-online}
\hlchaptermodality{204}

\begin{chapteropening}
\textbf{By the end of this chapter:} \emph{I can talk about websites, logging in, clicking, videos and typing.}

Complete the last lesson of chapter 204: I can talk about websites, logging in, clicking, videos and typing.
\end{chapteropening}

\section[\texorpdfstring{\zh{网站} (wǎngzhàn)}{网站 (wǎngzhàn)}]{\zh{网站} (wǎngzhàn) — website}

\label{lesson:ZH-C204-wangzhan}



\begin{quote}
Before the new one: say the Mandarin for a deposit paid in advance, then the Mandarin for contract.
\end{quote}

\subsection*{You'll want to know: \zh{网站}}
\textbf{\zh{网站}} — \emph{wǎngzhàn} — \textquotedblleft{}website\textquotedblright{}.

\begin{quote}
\textbf{\zh{这个网站关了}} \emph{zhège wǎngzhàn guān le} — This website has closed down.
\end{quote}

\begin{cousinweb}
\textbf{\zh{网站}} (\emph{wǎngzhàn}) is built from \zh{网} 'net' and \zh{站} 'a station', both inherited from Middle Chinese, \emph{mjangX} and \emph{treamH} (Guangyun, in Baxter's transcription).
\end{cousinweb}

\begin{grammarlens}[title={What you've built}]
One more word for reading notices and messages.
\end{grammarlens}

\subsection*{Your turn}
\begin{itemize}
\raggedright
  \item \emph{Say it:} \emph{wǎngzhàn}
  \item \emph{Say it:} \emph{wǎngzhàn}, once more
  \item \emph{Say it:} \emph{zhège wǎngzhàn guān le} — This website has closed down.
  \item \emph{Recall:} say the Mandarin for a deposit paid in advance, then the Mandarin for contract, then say \emph{wǎngzhàn} again
\end{itemize}

\begin{tcolorbox}[breakable,colback=teal!4,colframe=teal!35!black,title={Before you move on}]
What does \zh{网站} mean? (\textquotedblleft{}Website\textquotedblright{}.) Say it once more.
\end{tcolorbox}
\section[dēnglù]{dēnglù — to log in}

\label{lesson:ZH-C204-denglu}



\begin{quote}
Before the new one: say the Mandarin for contract, then the Mandarin for website.
\end{quote}

\subsection*{You'll want to know: dēnglù}
\emph{dēnglù} — \textquotedblleft{}to log in\textquotedblright{}.

\begin{quote}
\emph{nǐ dēnglù le ma} — Have you logged in?
\end{quote}

Its characters come later; for now it is a word for the ear.

\begin{cousinweb}
\emph{dēnglù} is built from \emph{dēng} 'to go up' and \emph{lù} 'to record', both inherited from Middle Chinese, \emph{tong} and \emph{ljowk} (Guangyun, in Baxter's transcription), which glosses the first 'to climb' and the second 'to take down'. Logging in is a modern sense, and how the pair came to it is not settled.
\end{cousinweb}

\begin{grammarlens}[title={What you've built}]
One more word for reading notices and messages.
\end{grammarlens}

\subsection*{Your turn}
\begin{itemize}
\raggedright
  \item \emph{Say it:} \emph{dēnglù}
  \item \emph{Say it:} \emph{dēnglù}, once more
  \item \emph{Say it:} \emph{nǐ dēnglù le ma} — Have you logged in?
  \item \emph{Recall:} say the Mandarin for contract, then the Mandarin for website, then say \emph{dēnglù} again
\end{itemize}

\begin{tcolorbox}[breakable,colback=teal!4,colframe=teal!35!black,title={Before you move on}]
What does dēnglù mean? (\textquotedblleft{}To log in\textquotedblright{}.) Say it once more.
\end{tcolorbox}
\section[diǎnjī]{diǎnjī — to click (on a screen)}

\label{lesson:ZH-C204-dianji}



\begin{quote}
Before the new one: say the Mandarin for website, then the Mandarin for to log in.
\end{quote}

\subsection*{You'll want to know: diǎnjī}
\emph{diǎnjī} — \textquotedblleft{}to click (on a screen)\textquotedblright{}.

\begin{quote}
\emph{qǐng diǎnjī zhèr} — Please click here.
\end{quote}

Its characters come later; for now it is a word for the ear.

\begin{cousinweb}
\emph{diǎnjī} is built from \zh{点} 'a dot, to touch' and \emph{jī} 'to strike', both inherited from Middle Chinese, \emph{temX} and \emph{kek} (Guangyun, in Baxter's transcription), which glosses the second 'to hit'.
\end{cousinweb}

\begin{grammarlens}[title={What you've built}]
One more word for reading notices and messages.
\end{grammarlens}

\subsection*{Your turn}
\begin{itemize}
\raggedright
  \item \emph{Say it:} \emph{diǎnjī}
  \item \emph{Say it:} \emph{diǎnjī}, once more
  \item \emph{Say it:} \emph{qǐng diǎnjī zhèr} — Please click here.
  \item \emph{Recall:} say the Mandarin for website, then the Mandarin for to log in, then say \emph{diǎnjī} again
\end{itemize}

\begin{tcolorbox}[breakable,colback=teal!4,colframe=teal!35!black,title={Before you move on}]
What does diǎnjī mean? (\textquotedblleft{}To click (on a screen)\textquotedblright{}.) Say it once more.
\end{tcolorbox}
\section[shìpín]{shìpín — video}

\label{lesson:ZH-C204-shipin}



\begin{quote}
Before the new one: say the Mandarin for to log in, then the Mandarin for to click.
\end{quote}

\subsection*{You'll want to know: shìpín}
\emph{shìpín} — \textquotedblleft{}video\textquotedblright{}.

\begin{quote}
\emph{wǒ xǐhuan kàn shìpín} — I like watching videos.
\end{quote}

Its characters come later; for now it is a word for the ear.

\begin{cousinweb}
\emph{shìpín} is built from \emph{shì} 'to look at' and \emph{pín} 'frequent', both inherited from Middle Chinese, \emph{dzyijH} and \emph{bjin} (Guangyun, in Baxter's transcription). The word is modern, and how the pair came to mean a video is not settled.
\end{cousinweb}

\begin{grammarlens}[title={What you've built}]
One more word for reading notices and messages.
\end{grammarlens}

\subsection*{Your turn}
\begin{itemize}
\raggedright
  \item \emph{Say it:} \emph{shìpín}
  \item \emph{Say it:} \emph{shìpín}, once more
  \item \emph{Say it:} \emph{wǒ xǐhuan kàn shìpín} — I like watching videos.
  \item \emph{Recall:} say the Mandarin for to log in, then the Mandarin for to click, then say \emph{shìpín} again
\end{itemize}

\begin{tcolorbox}[breakable,colback=teal!4,colframe=teal!35!black,title={Before you move on}]
What does shìpín mean? (\textquotedblleft{}Video\textquotedblright{}.) Say it once more.
\end{tcolorbox}
\section[\texorpdfstring{\zh{打字} (dǎzì)}{打字 (dǎzì)}]{\zh{打字} (dǎzì) — to type}

\label{lesson:ZH-C204-dazi}



\begin{quote}
Before the new one: say the Mandarin for to click, then the Mandarin for video.
\end{quote}

\subsection*{You'll want to know: \zh{打字}}
\textbf{\zh{打字}} — \emph{dǎzì} — \textquotedblleft{}to type\textquotedblright{}.

\begin{quote}
\textbf{\zh{他打字很快}} \emph{tā dǎzì hěn kuài} — He types fast.
\end{quote}

\begin{cousinweb}
\textbf{\zh{打字}} (\emph{dǎzì}) is built from \zh{打} 'to strike' and \zh{字} 'a character', both inherited from Middle Chinese, \emph{taengX} and \emph{dziH} (Guangyun, in Baxter's transcription): striking out characters. The modern dǎ is not the regular outcome of \emph{taengX}, and how it came is not settled.
\end{cousinweb}

\begin{grammarlens}[title={What you've built}]
One more word for reading notices and messages.
\end{grammarlens}

\subsection*{Your turn}
\begin{itemize}
\raggedright
  \item \emph{Say it:} \emph{dǎzì}
  \item \emph{Say it:} \emph{dǎzì}, once more
  \item \emph{Say it:} \emph{tā dǎzì hěn kuài} — He types fast.
  \item \emph{Recall:} say the Mandarin for to click, then the Mandarin for video, then say \emph{dǎzì} again
\end{itemize}

\begin{tcolorbox}[breakable,colback=teal!4,colframe=teal!35!black,title={Before you move on}]
What does \zh{打字} mean? (\textquotedblleft{}To type\textquotedblright{}.) Say it once more.
\end{tcolorbox}
